\documentclass[11pt,a4paper]{article}
\usepackage[margin=25mm]{geometry}
\usepackage{amsmath,amssymb,array,graphicx,hyperref}
\hypersetup{colorlinks=true,urlcolor=blue,linkcolor=blue,pdftitle={Rigid-body contact in combined velocity and momentum notation},pdfauthor={Viktor: notation and model}}
\newcommand{\rw}[1]{#1\mathbin{\wedge}}
\newcommand{\lwedge}[1]{\mathbin{\wedge}#1}
\newcommand{\id}{\mathbf1}
\newcommand{\inertia}{\mathbb I}
\title{Rigid-body contact in combined velocity and momentum notation\\\large Symbolic mechanics, directional impulses and coupled matrix calculations}
\author{Viktor's notation: $\mathbf r\wedge$, $\ell\boldsymbol\omega$ and $\mathbf L/\ell$}
\date{6 October 2026}
\begin{document}
\maketitle
\begin{abstract}
We combine linear velocity and length-scaled angular velocity in a single vector, and combine linear momentum and reciprocally scaled angular momentum in its dual. A baseline wedge denotes the cross-product matrix; moving the wedge to the other side denotes its transpose. The resulting additive matrix chains include both the moment of the linear impulse and an independent angular impulse. We derive the contact response, its two-body and simultaneous-contact assembly, kinetic-energy balance, and matrix-free evaluation in two and three dimensions. Sliding and spin directions follow the relative motions rather than an arbitrary contact basis. Numerical results are omitted pending recalculation with the intended directional contact law.
\end{abstract}

\section{Notation and scope}
Lowercase $\mathbf p$ is linear momentum; uppercase $\mathbf L$ is angular momentum. Lowercase $\delta$ denotes an impulse, and component indices are lowercase. Uppercase $\mathbf V$, $\mathbf P$ and $\mathbf F$ denote combined motion, momentum and force vectors. The fixed reference length is $\ell$; it is not a material coefficient.

The direction $\mathbf t$ follows relative velocity at the collision point, and $\mathbf s$ follows relative angular velocity. They are motion-defined directions, not two axes chosen to complete a tangent basis. The contact normal is $\mathbf n$. The allowed directional impulse components must be respected; a linear spin component $\delta\mathbf p_s$ is not introduced.

The directional impulse labels are $\delta\mathbf p_n$, $\delta\mathbf p_t$, $\delta\mathbf L_s$ and $\delta\mathbf L_n$. The normal and sliding linear impulses are distinct from the spin and normal angular impulses.

The mechanics derived here is exact for a configuration frozen during an instantaneous impulse. It maps an allowed contact impulse to motion; it does not, by itself, determine the impulse. The original directional restitution and friction equations must close that calculation. Where those equations have not been recovered, they remain symbolic rather than being replaced by a conventional force-only law.

\section{A baseline wedge for the cross-product matrix}
In three dimensions define
\begin{equation}
\rw{\mathbf r}=
\begin{bmatrix}
0&-r_z&r_y\\r_z&0&-r_x\\-r_y&r_x&0
\end{bmatrix},\qquad
\rw{\mathbf r}\mathbf a=\mathbf r\times\mathbf a.
\end{equation}
The wedge sits on the same baseline as the vector. The reversed notation denotes the transpose:
\begin{equation}
\lwedge{\mathbf r}=\rw{\mathbf r}^{\mathsf T}.
\end{equation}
Here the superscript $\mathsf T$ applies to the matrix represented by $\mathbf r\wedge$. Moving the wedge from right to left denotes the transpose in both dimensions. No minus-sign replacement or operand-reversal rule is used.

Matrix associativity permits chains without redundant parentheses, for example
\begin{equation}
\lwedge{\mathbf r}\inertia^{-1}\rw{\mathbf r}\delta\mathbf p.
\end{equation}
The factor order is retained. Transposition reverses a product's factor order; it does not make general matrices commute.

In two dimensions the moment map is a row and its transpose is a column:
\begin{equation}
\rw{\mathbf r}=\begin{bmatrix}-r_y&r_x\end{bmatrix},\qquad
\lwedge{\mathbf r}=\begin{bmatrix}-r_y\\r_x\end{bmatrix}.
\end{equation}
Thus $\rw{\mathbf r}\mathbf a=r_xa_y-r_ya_x$. The right-wedge map takes a planar linear vector to a scalar moment, while its left-wedge transpose takes a scalar angular velocity to a planar contact-velocity contribution. The same transpose notation therefore applies to rectangular planar maps and square spatial maps.

\section{Combined body velocity, momentum and force}
Let $m$ be mass, $\inertia$ the world-frame central inertia, $\mathbf v$ the centre-of-mass velocity and $\boldsymbol\omega$ angular velocity. Define
\begin{equation}
\mathbf V=\begin{bmatrix}\mathbf v\\\ell\boldsymbol\omega\end{bmatrix},\quad
\mathbf P=\begin{bmatrix}\mathbf p\\\mathbf L/\ell\end{bmatrix},\quad
\mathbf M=\begin{bmatrix}m\id&0\\0&\inertia/\ell^2\end{bmatrix},\quad
\boxed{\mathbf P=\mathbf M\mathbf V}.
\end{equation}
The combined force and its impulse are
\begin{equation}
\mathbf F=\begin{bmatrix}\mathbf f\\\boldsymbol\tau/\ell\end{bmatrix},\qquad
\delta\mathbf P=\int\mathbf F\,dt,
\end{equation}
where $\boldsymbol\tau$ is the torque about the centre of mass for this body-level expression. Work and kinetic energy retain their physical values:
\begin{equation}
\mathbf V^{\mathsf T}\mathbf F
=\mathbf v^{\mathsf T}\mathbf f+\boldsymbol\omega^{\mathsf T}\boldsymbol\tau,\qquad
E_k=\tfrac12\mathbf V^{\mathsf T}\mathbf M\mathbf V.
\end{equation}
In planar motion the vectors have three components, the angular quantities are scalars, and $\mathbf M=\operatorname{diag}(m,m,\inertia/\ell^2)$.

A change of fixed representation length must leave the physical velocity, impulse and energy unchanged. Changing $\ell$ during motion would require coordinate-derivative terms and is outside this presentation. For simplicity we use one common fixed length throughout the contact assembly.

\section{Contact motion and the full impulse}
Let the contact reference point be $\mathbf x_c$ and the body centre be $\mathbf x_a$. Its lever arm is $\mathbf r_a=\mathbf x_c-\mathbf x_a$. The contact velocity is additive:
\begin{equation}
\mathbf v_{c,a}=\mathbf v_a+\boldsymbol\omega_a\times\mathbf r_a
=\mathbf v_a+\lwedge{\mathbf r_a}\boldsymbol\omega_a.
\end{equation}
Set
\begin{equation}
\mathbf C_a=\begin{bmatrix}\id&\lwedge{\mathbf r_a}/\ell\end{bmatrix},\qquad
\mathbf B_a=\begin{bmatrix}\id&\lwedge{\mathbf r_a}/\ell\\0&\id\end{bmatrix}.
\end{equation}
Then the full contact motion is
\begin{equation}
\mathbf V_{c,a}=\begin{bmatrix}\mathbf v_{c,a}\\\ell\boldsymbol\omega_a\end{bmatrix}
=\mathbf B_a\mathbf V_a.
\end{equation}
At the declared contact reference, define the linear impulse $\delta\mathbf p$ and the independent angular impulse $\delta\mathbf L$:
\begin{equation}
\delta\mathbf p=\int\mathbf f_c\,dt,\qquad
\delta\mathbf L=\int\boldsymbol\tau_c\,dt,\qquad
\delta\mathbf P_c=\begin{bmatrix}\delta\mathbf p\\\delta\mathbf L/\ell\end{bmatrix}.
\end{equation}
The subscript $c$ distinguishes this contact impulse from the body momentum increment. The independent contact angular impulse excludes the lever-arm moment, which the transpose map supplies:
\begin{equation}
\boxed{\delta\mathbf P_a=\mathbf B_a^{\mathsf T}\delta\mathbf P_c}
=\begin{bmatrix}\delta\mathbf p\\
\rw{\mathbf r_a}\delta\mathbf p/\ell+\delta\mathbf L/\ell\end{bmatrix}.
\end{equation}
Equivalently,
\begin{equation}
\delta\mathbf p_a=\delta\mathbf p,\qquad
\boxed{\delta\mathbf L_a=\rw{\mathbf r_a}\delta\mathbf p+\delta\mathbf L}.
\end{equation}
Here $\delta\mathbf L_a$ is the total body angular impulse. This distinction prevents counting the force moment twice. The associated work is
\begin{equation}
\mathbf V_a^{\mathsf T}\delta\mathbf P_a
=\mathbf V_{c,a}^{\mathsf T}\delta\mathbf P_c
=\mathbf v_{c,a}^{\mathsf T}\delta\mathbf p
+\boldsymbol\omega_a^{\mathsf T}\delta\mathbf L.
\end{equation}
For distributed contact traction $\mathbf f_{\!A}$ the independent angular impulse is the traction moment about the chosen reference:
\begin{equation}
\delta\mathbf p=\int\!\int_{\mathcal A}\mathbf f_{\!A}\,dA\,dt,\qquad
\delta\mathbf L=\int\!\int_{\mathcal A}\rw{\mathbf x-\mathbf x_c}\mathbf f_{\!A}\,dA\,dt.
\end{equation}
A reference shift $\mathbf x_c'=\mathbf x_c+\mathbf a$ gives
\begin{equation}
\mathbf r_a'=\mathbf r_a+\mathbf a,\qquad
\delta\mathbf L'=\delta\mathbf L-\rw{\mathbf a}\delta\mathbf p,
\end{equation}
leaving the total body angular impulse unchanged. This is reference bookkeeping, not a way to fit a collision outcome by moving the force point.

\section{Two-body relative motion and contact mobility}
Orient impulse transfer from body B to body A, at one shared reference point. Body A receives $+\delta\mathbf P_c$ and body B receives $-\delta\mathbf P_c$. Define
\begin{equation}
\mathbf u=\mathbf v_{c,A}-\mathbf v_{c,B},\qquad
\boldsymbol\Omega=\boldsymbol\omega_A-\boldsymbol\omega_B,\qquad
\mathbf U=\begin{bmatrix}\mathbf u\\\ell\boldsymbol\Omega\end{bmatrix}.
\end{equation}
For stacked body variables and their mass matrix,
\begin{equation}
\mathbf G=\begin{bmatrix}\mathbf B_A&-\mathbf B_B\end{bmatrix},\quad
\mathbf M=\operatorname{diag}(\mathbf M_A,\mathbf M_B),\quad
\mathbf U=\mathbf G\mathbf V.
\end{equation}
The impulse and motion updates are
\begin{align}
\delta\mathbf P&=\mathbf G^{\mathsf T}\delta\mathbf P_c,\\
\mathbf V^+&=\mathbf V^-+\mathbf M^{-1}\mathbf G^{\mathsf T}\delta\mathbf P_c,\\
\boxed{\mathbf U^+=\mathbf U^-+\mathbf W\delta\mathbf P_c},\qquad
\mathbf W&=\mathbf G\mathbf M^{-1}\mathbf G^{\mathsf T}.
\end{align}
The explicit two-body matrix is
\begin{equation}
\boxed{\mathbf W=\sum_{a\in\{A,B\}}
\begin{bmatrix}
\id/m_a+\lwedge{\mathbf r_a}\inertia_a^{-1}\rw{\mathbf r_a}
&\ell\lwedge{\mathbf r_a}\inertia_a^{-1}\\
\ell\inertia_a^{-1}\rw{\mathbf r_a}&\ell^2\inertia_a^{-1}
\end{bmatrix}}.
\end{equation}
All products retain their order and the off-diagonal blocks couple linear and angular impulses. The matrix is symmetric positive semidefinite. A prescribed boundary contributes its known motion but zero inverse mass; its mass is not inserted as a large finite substitute.

\section{Sliding and spin directions}
For nonzero relative motions the definitions supplied by the model are
\begin{equation}
\boxed{\mathbf t=\mathbf u/\|\mathbf u\|},\qquad
\boxed{\mathbf s=\boldsymbol\Omega/\|\boldsymbol\Omega\|}.
\end{equation}
These definitions use relative velocity \emph{at the collision point}, not centre-of-mass velocity. No projection onto a geometric tangent plane has been inserted. Consequently $\mathbf t$ need not be perpendicular to $\mathbf n$, and $\mathbf s$ need not be perpendicular to either. If the original law specifies a projected sliding velocity, that projection must be recovered explicitly rather than assumed.

Write the allowed normal/sliding linear and spin/normal angular components as
\begin{align}
\delta\mathbf p_n=q_n\mathbf n,\qquad
\delta\mathbf p_t=q_t\mathbf t,\qquad
\delta\mathbf p&=\delta\mathbf p_n+\delta\mathbf p_t,\\
\delta\mathbf L_s=a_s\mathbf s,\qquad
\delta\mathbf L_n=a_n\mathbf n,\qquad
\delta\mathbf L&=\delta\mathbf L_s+\delta\mathbf L_n.
\end{align}
The scalars $q_n,q_t$ specify linear impulse amplitudes and $a_s,a_n$ specify angular impulse amplitudes. Introduce the component map
\begin{equation}
\mathbf D=\begin{bmatrix}
\mathbf n&\mathbf t&0&0\\
0&0&\mathbf s&\mathbf n
\end{bmatrix},\qquad
\boldsymbol\lambda=\begin{bmatrix}q_n\\q_t\\a_s/\ell\\a_n/\ell\end{bmatrix},\qquad
\boxed{\delta\mathbf P_c=\mathbf D\boldsymbol\lambda}.
\end{equation}
This map is not assumed orthogonal, square or invertible. Its dual gives the component-conjugate motions and coupled mobility:
\begin{equation}
\mathbf z=\mathbf D^{\mathsf T}\mathbf U,\qquad
\mathbf A=\mathbf D^{\mathsf T}\mathbf W\mathbf D,\qquad
\boxed{\mathbf z^+=\mathbf z^-+\mathbf A\boldsymbol\lambda}.
\end{equation}
The transpose is a work projection, not an assertion that nonorthogonal directions are independent coordinate axes.

If a relative motion vanishes, its unit direction is undefined. Its subsequent static, sticking or onset branch must come from the contact law. Choosing an arbitrary direction and calling it sliding or spin changes the model. If $\mathbf s$ aligns with $\mathbf n$, the two angular columns become dependent. The law must then determine their physical sum without demanding a nonsingular four-component inverse. The same issue can occur for the linear directions.

Directions evaluated before a frozen impulse give a finite matrix calculation. Directions evolving during a finite collision require state updates and force integration. Those are distinct algorithms and must not be conflated.

\section{Restitution and friction as the contact closure}
The mass and contact matrices specify how impulses change motion. Restitution and friction specify which impulses are applied. Keep their material parameters in a separate profile $\boldsymbol\theta$ and express the directional law symbolically as
\begin{equation}
\boldsymbol\Phi\bigl(\boldsymbol\lambda,\mathbf U^-,\mathbf U^+,\mathbf D,
\boldsymbol\theta,\boldsymbol\xi\bigr)=0,\qquad
\boldsymbol\lambda\in\mathcal A\bigl(\boldsymbol\theta,\boldsymbol\xi\bigr).
\end{equation}
Here $\boldsymbol\xi$ denotes contact history only if the original law retains it, and $\mathcal A$ denotes its physically specified admissibility conditions. These placeholders identify the missing constitutive specification; they are not a proposed new fitted law.

For any scalar motion channel to which a signed endpoint restitution condition actually applies, its symbolic target is
\begin{equation}
z_i^++e_i z_i^-=0.
\end{equation}
For a declared set of such channels, a frozen-direction residual has the form
\begin{equation}
\mathbf A\boldsymbol\lambda+\id\mathbf z^-+\mathbf E\mathbf z^-=0,
\end{equation}
with $\mathbf E$ containing the corresponding restitution symbols. This is a conditional matrix expression, not a statement that all four components have independently specified restitution targets. Friction limits, dependent directions or simultaneous-contact kinematics can prevent exact targets from coexisting.

Normal restitution and the historical geometric-tangent restitution must not be silently applied to a different motion-defined direction. A usual signed restitution convention means $e_i=-1$ preserves the channel velocity, $e_i=0$ targets arrest, and $e_i>0$ targets reversal. That convention alone establishes neither a material property nor passive coupled motion.

Sliding and angular resistance must follow the model's specified $\mathbf t$, $\mathbf s$ and normal angular component. A translational Coulomb cone by itself supplies no independent angular impulse. Conversely, adding independent guessed moment limits or replacing the direction map by two arbitrary tangent axes changes the constitutive model. We therefore do not retain the earlier force-only endpoint solver as a definition of this law.

\begin{center}\small
\begin{tabular}{p{40mm}p{93mm}}\hline
Separate input field & Required specification\\\hline
Normal restitution & Documented material-pair value and velocity convention\\
Sliding restitution & Documented value and the exact relative-motion channel\\
Sliding friction & Original static/dynamic or endpoint rule, if present\\
Angular contact response & Original $s$ and $n$ component laws and their parameters\\
Geometry and state & Contact reference, normal, inertia, pose, incoming motion\\
Validity & Surface condition, speed/size range, uncertainties and source\\\hline
\end{tabular}\end{center}
No numerical material value is substituted into the symbolic derivation. Missing public inputs remain missing. No coefficients are fitted to the target collision to manufacture agreement.

\section{Energy, angular momentum and moving boundaries}
For two dynamic bodies with no external impulsive work, the frozen-configuration kinetic-energy change is
\begin{equation}
\boxed{\delta E_k=\delta\mathbf P_c^{\mathsf T}\mathbf U^-
+\tfrac12\delta\mathbf P_c^{\mathsf T}\mathbf W\delta\mathbf P_c}.
\end{equation}
Its linear term contains \emph{both} impulse works:
\begin{equation}
\delta\mathbf P_c^{\mathsf T}\mathbf U^-
=\delta\mathbf p^{\mathsf T}\mathbf u^-
+\delta\mathbf L^{\mathsf T}\boldsymbol\Omega^-.
\end{equation}
In the allowed component representation,
\begin{equation}
\boxed{\delta E_k=\boldsymbol\lambda^{\mathsf T}\mathbf z^-
+\tfrac12\boldsymbol\lambda^{\mathsf T}\mathbf A\boldsymbol\lambda}.
\end{equation}
For an initially unstrained passive impact require $\delta E_k\le0$. If stored elastic energy is present, include it in the ledger instead of imposing a kinetic-energy inequality in isolation. Translation alone need not conserve energy, because rotation, deformation, dissipation and support work exchange energy.

For a prescribed moving boundary the work supplied through its contact motion is
\begin{equation}
W_b=\delta\mathbf p^{\mathsf T}\mathbf v_{c,b}
+\delta\mathbf L^{\mathsf T}\boldsymbol\omega_b,
\end{equation}
and the corresponding relative-motion identity concerns $\delta E_k-W_b$.

Equal and opposite impulses at the shared reference conserve total linear momentum. About a world origin the angular balance uses the orbital contribution as well as central angular momentum:
\begin{equation}
\delta\mathbf J_a=\rw{\mathbf x_a}\delta\mathbf p_a+\delta\mathbf L_a.
\end{equation}
The pair contributions sum to zero, including the independent opposite angular impulses. Checking only the force lever-arm moment misses this mechanism.

\section{Symbolic three-dimensional examples}
\subsection{Central normal impact with axial spin}
Consider a homogeneous sphere of radius $R$ and central inertia $\inertia\id$, with $\mathbf r=-R\mathbf n$. Prescribe a central linear impulse and axial incoming spin:
\begin{equation}
\delta\mathbf p=q\mathbf n,\qquad
\boldsymbol\omega^- =\omega^-\mathbf n.
\end{equation}
Then
\begin{equation}
\rw{\mathbf r}\delta\mathbf p=0,\qquad
\lwedge{\mathbf r}\boldsymbol\omega^-=0.
\end{equation}
The complete angular update is therefore
\begin{equation}
\boxed{\delta\mathbf L_a=\delta\mathbf L_s+\delta\mathbf L_n},\qquad
\boldsymbol\omega^+=\boldsymbol\omega^-+\inertia^{-1}\delta\mathbf L_a.
\end{equation}
A force-only model sets these independent angular impulses to zero and cannot change this axial spin. Here $\mathbf s$ is parallel or antiparallel to $\mathbf n$: the component map is dependent, yet the physical total angular impulse is well defined. The contact law must determine that total. No arbitrary inversion of two collinear angular directions is required.

\subsection{Rotated non-spherical body}
Let a homogeneous cuboid have half-extents $a,b,c$, mass $m$, and attitude matrix $\mathbf Q$. Its shape-frame and world inertia are
\begin{equation}
\inertia_0=\frac m3\operatorname{diag}(b^2+c^2,a^2+c^2,a^2+b^2),\qquad
\inertia=\mathbf Q\inertia_0\mathbf Q^{\mathsf T}.
\end{equation}
For a supporting vertex $\mathbf r_0$, use $\mathbf r=\mathbf Q\mathbf r_0$. A fixed plane gives
\begin{equation}
\mathbf u^- =\mathbf v^-+\lwedge{\mathbf r}\boldsymbol\omega^-,\qquad
\boldsymbol\Omega^- =\boldsymbol\omega^-.
\end{equation}
Construct $\mathbf t$, $\mathbf s$, the direction map and the full mobility from these physical states. Once the directional law supplies $\delta\mathbf p$ and $\delta\mathbf L$, predict
\begin{equation}
\boxed{\mathbf v^+=\mathbf v^-+\delta\mathbf p/m},\qquad
\boxed{\boldsymbol\omega^+=\boldsymbol\omega^-
+\inertia^{-1}\rw{\mathbf r}\delta\mathbf p
+\inertia^{-1}\delta\mathbf L}.
\end{equation}
The same formula includes spin-generated sliding, an off-centre normal force moment, and the independent angular impulse. Non-diagonal inertia couples the angular directions without adding material parameters. This is a symbolic 3D example; the earlier force-only cuboid numbers have been removed from this reading version.

\section{The planar form}
For a planar contact set
\begin{equation}
\mathbf V_a=\begin{bmatrix}v_{x,a}\\v_{y,a}\\\ell\omega_a\end{bmatrix},\qquad
\mathbf B_a=\begin{bmatrix}1&0&-r_{y,a}/\ell\\0&1&r_{x,a}/\ell\\0&0&1\end{bmatrix}.
\end{equation}
The contact and body impulses are
\begin{equation}
\delta\mathbf P_c=\begin{bmatrix}\delta p_x\\\delta p_y\\\delta L/\ell\end{bmatrix},\qquad
\delta\mathbf P_a=\mathbf B_a^{\mathsf T}\delta\mathbf P_c,\qquad
\delta L_a=\rw{\mathbf r_a}\delta\mathbf p+\delta L.
\end{equation}
The pair map, mobility and energy identities retain the same matrix forms with these planar block dimensions. Planar angular velocity and angular impulse are signed scalars. A three-dimensional normal-angular channel is restricted to the permitted planar rotation axis; it must not be inserted as an extra independent planar degree of freedom. Any component unavailable under that geometric restriction is omitted rather than independently solved.

\section{Simultaneous contacts and efficient evaluation}
For many bodies, stack their combined velocities and block-diagonal masses. For each oriented contact, put $+\mathbf B_a$ and $-\mathbf B_b$ in the appropriate columns of its row block. Stack the allowed component maps into $\mathcal D$. Define
\begin{equation}
\mathbf H=\mathcal D^{\mathsf T}\mathbf G,\qquad
\mathbf z=\mathbf H\mathbf V,\qquad
\mathbf A=\mathbf H\mathbf M^{-1}\mathbf H^{\mathsf T}.
\end{equation}
The whole contact system is
\begin{equation}
\mathbf z^+=\mathbf z^-+\mathbf A\boldsymbol\lambda,\qquad
\mathbf V^+=\mathbf V^-+\mathbf M^{-1}\mathbf H^{\mathsf T}\boldsymbol\lambda.
\end{equation}
Off-diagonal contact blocks carry impulses through shared bodies. Keeping only independent pair blocks changes rapid group motion. Disconnected contact islands may be solved separately.

The operator need not be formed densely. For a trial component vector $\mathbf q$, apply
\begin{equation}
\boxed{\mathbf A\mathbf q=\mathbf H\mathbf M^{-1}\mathbf H^{\mathsf T}\mathbf q}
\end{equation}
through three operations: scatter $\mathbf H^{\mathsf T}\mathbf q$ to bodies, apply each body's inverse mass and inertia solve, and gather with $\mathbf H$. Rotation uses the world inertia $\mathbf Q\inertia_0\mathbf Q^{\mathsf T}$. Solve small inertia systems or use a rotated precomputed shape-frame inverse; avoid forming a large inverse.

The directional law selects the component unknowns and active branches. Assemble small matrices for modest islands or use iterative operator evaluations for large islands. Dependent directions and redundant contacts require rank-aware solves; a failed inverse is not repaired by silently changing the physical components. Check the full-system work and energy, not independent contact energies that omit shared-body cross terms.

\section{Recalculation and comparison plan}
Numerical outputs, material-value tables and prediction-error diagrams are intentionally omitted from this rewrite. Existing results remain preserved in the repository as calculations of historical comparator models. They are not results of the directional model described here.

Before new predictions, recover the exact original restitution/friction closure, its zero-motion branches, and whether directions evolve during impact. Then verify pure spin, pure sliding, mixed motion, off-centre non-spherical impacts and coupled contact groups in both dimensions. Checks must include equal/opposite impulse, total angular momentum, dual work, energy and reference-length invariance.

Select independently documented material-pair parameters without project fitting. Missing angular response data cannot be inferred from desired output and called an independent material prediction. Public ball and rock records can then be compared using the same measured input states, with actual shape, inertia and uncertainty wherever available. Source reconstruction and held-out prediction must remain distinguishable.

For the eventual visualization, compare measured and predicted outgoing motion using the previously requested two-coordinate deviation map. Define the observation signature and its metric before inspecting agreement; state which components were independently measured. Distinct collision types can share the same error point. No dots are supplied here until the corrected model has actually been evaluated.

\section{Relation to prior notation}
Combined motion and dual force/momentum, power pairing and characteristic-length normalization have established precedents\cite{featherstone,modernrobotics,angeles}. Vose and colleagues explicitly scale planar angular velocity and torque by reciprocal factors\cite{vose}; Zhang and colleagues give related deformation/moment scaling\cite{zhang}. The baseline wedge here makes the cross-product transpose and additive matrix chains explicit; matrix associativity and transpose duality are established algebra.

Contact force-and-moment models and matrix-free coupled contact calculations also have prior art\cite{drakewrench,mujocowrench,siggraph}. This symbolic rewrite is a presentation of the intended notation and mechanics, not a demonstrated novelty claim or a completed validation of the directional constitutive law. The separate literature review retains the source-specific findings. A publishable predictive claim requires new, reproducible evidence after the corrected model is implemented.

\begingroup
\raggedright
\begin{thebibliography}{20}
\bibitem{cornell} Cornell Grain Flow Research, impact parameter chart and experimental method. \url{https://grainflowresearch.mae.cornell.edu/impact/data/Impact%20Results.html}; \url{https://grainflowresearch.mae.cornell.edu/impact/impact.html}.
\bibitem{joseph} G. G. Joseph and M. L. Hunt, \emph{Oblique particle-wall collisions in a liquid}, J. Fluid Mech. 510 (2004), 71--93. Dry profiles on p78. \url{https://doi.org/10.1017/S002211200400919X}.
\bibitem{cross} R. Cross, \emph{Enhancing the Bounce of a Ball} (2010), Table I. Author-hosted manuscript: \url{https://physics.usyd.edu.au/~cross/PUBLICATIONS/48.%20EnhanceBounce.pdf}.
\bibitem{wang} Wang et al., \emph{Effects of the impact angle on the coefficient of restitution in rockfall}, Natural Hazards and Earth System Sciences 18 (2018), 3045--3060. \url{https://doi.org/10.5194/nhess-18-3045-2018}. Publisher supplementary spreadsheet provides the impact observations.
\bibitem{featherstone} R. Featherstone, \emph{General Comments on Coordinate Transforms}, Spatial Vector and Rigid-Body Dynamics documentation. Motion/force duality and inertia transformation. \url{https://royfeatherstone.org/spatial/v2/xforms.html}.
\bibitem{modernrobotics} K. M. Lynch and F. C. Park, \emph{Modern Robotics}, section 3.4, Wrenches. Twist--wrench power pairing. \url{https://modernrobotics.northwestern.edu/nu-gm-book-resource/3-4-wrenches/}.
\bibitem{angeles} J. Angeles, \emph{On the Optimum Dimensioning of Robotic Manipulators}, author-hosted manuscript, sections 1--2. Characteristic-length normalization and homogeneous Jacobians. \url{https://cim.mcgill.ca/~rmsl/Index/Documents/MECH573/robotdes0312.pdf}.
\bibitem{vose} T. H. Vose, P. Umbanhowar and K. M. Lynch, \emph{Sliding Manipulation of Rigid Bodies on a Controlled 6-DoF Plate}, Robotics: Science and Systems VII (2011). Appendix footnote 1, PDF p7. \url{https://www.roboticsproceedings.org/rss07/p43.pdf}.
\bibitem{zhang} Y. Zhang, H.-J. Su and Q. Liao, \emph{Mobility criteria of compliant mechanisms based on decomposition of compliance matrices}, Mechanism and Machine Theory 79 (2014), 80--93. Section 3.3, equation 21. \url{https://doi.org/10.1016/j.mechmachtheory.2014.04.010}.
\bibitem{siggraph} \emph{Contact and Friction Simulation for Computer Graphics}, SIGGRAPH 2022 course notes. Matrix-free tutorial, PDF pp163--164. \url{https://siggraphcontact.github.io/assets/files/SIGGRAPH22_friction_contact_notes.pdf}.
\bibitem{acary} V. Acary and N. A. Collins-Craft, \emph{On the Moreau--Jean scheme with the Fr\'emond impact law: energy conservation and dissipation properties for elastodynamics with contact, impact and friction}, JTCAM (2025). Section 2.5. \url{https://doi.org/10.46298/jtcam.13480}.
\bibitem{drakewrench} Drake documentation, \emph{Hydroelastic Contact User Guide}. Distributed contact force and torque. \url{https://drake.mit.edu/doxygen_cxx/group__hydroelastic__user__guide.html}.
\bibitem{mujocowrench} MuJoCo documentation, \emph{Computation: Contact}. Torsional and rolling friction and higher-dimensional contact wrenches. \url{https://mujoco.readthedocs.io/en/stable/computation.html}.
\end{thebibliography}
\endgroup
\end{document}
